% Sezione 1: Introduzione e Contesto Aziendale
\section{Introduzione}

% Slide 1: Contesto Aziendale
\begin{frame}
    \frametitle{Contesto Aziendale: Daikin Applied Europe}
    \begin{columns}
        \column{0.6\textwidth}
        \begin{itemize}
            \item \textbf{Daikin Applied Europe (DAE)}:
                  \begin{itemize}
                      \item Sede di Ariccia (Roma).
                  \end{itemize}
            \item Leader nel settore HVAC (Heating, Ventilation, Air Conditioning).
            \item Produzione \textbf{Make-to-Order}:
                  \begin{itemize}
                      \item Unità Trattamento Aria (UTA) e Chiller.
                      \item Prodotti personalizzati per il mercato europeo.
                  \end{itemize}
            \item \textbf{Driver}: Crescente domanda per soluzioni in ambiti industriali rilevanti nel momento storico attuale (es: datacenter).
        \end{itemize}
        \column{0.4\textwidth}
        % Placeholder per logo o immagine stabilimento se disponibile
        \centering
        \includegraphics[width=\textwidth]{daikin/daikin_factory_ext.jpg} % Placeholder
    \end{columns}
\end{frame}

% Slide 2: Il Reparto Industrial Innovation
\begin{frame}
    \frametitle{Il Reparto Industrial Innovation}
    Obiettivo: Ottimizzazione dei processi produttivi verso l'\textbf{Industria 5.0}.
    \vspace{0.5cm}
    \begin{block}{Progetti Chiave}
        \begin{itemize}
            \item \textbf{MES (Manufacturing Execution System)}: Monitoraggio produzione in tempo reale.
            \item \textbf{Test-bench automatici}: Collaudo sottoassiemi e prodotti finiti.
            \item \textbf{Visione Artificiale}: Controllo qualità e ispezione.
            \item \textbf{Robotica Mobile}: Integrazione robot per movimentazione materiali.
        \end{itemize}
    \end{block}
\end{frame}

% Slide 3: Motivazioni e Problematiche
\begin{frame}
    \frametitle{Motivazioni e Problematiche}
    Attualmente la movimentazione è manuale (muletti, carrelli), con costi elevati in termini di tempo e risorse.
    \vspace{0.5cm}
    \begin{alertblock}{Sfide della Produzione Make-to-Order}
        \begin{itemize}
            \item \textbf{Percorsi Variabili}: Destinazioni cambiano in base alle esigenze giornaliere.
            \item \textbf{Layout Dinamico}: Disposizione spazi variabile per necessità logistiche.
            \item \textbf{Ostacoli Dinamici}: Presenza costante di operatori e altri mezzi.
        \end{itemize}
    \end{alertblock}
\end{frame}

% Slide 4: Obiettivo della Tesi
\begin{frame}
    \frametitle{Obiettivo della Tesi}
    Progettazione e realizzazione di un \textbf{Veicolo Autonomo (AMR)} per lo stabilimento di Ariccia.

    \begin{itemize}
        \item \textbf{Dimostratore Tecnologico}:
              \begin{itemize}
                  \item Realizzazione di un prototipo a \textbf{basso costo}.
                  \item Valutazione fattibilità in ambiente industriale reale, all'aperto e poco strutturato.
              \end{itemize}
        \item \textbf{Transizione all'Automazione}:
              \begin{itemize}
                  \item Sostituzione movimentazione manuale.
                  \item Integrazione futura con sistemi \textbf{WMS} (Warehouse Management System) ed \textbf{ERP}.
              \end{itemize}
    \end{itemize}
\end{frame}


